\documentclass[11pt]{article}
\usepackage[a4paper,margin=1in]{geometry}
\usepackage{amsmath,amssymb,amsthm}
\usepackage[T1]{fontenc}
\usepackage{lmodern}
\usepackage{microtype}
\usepackage{enumitem}
\usepackage{booktabs}
\usepackage{array}
\usepackage{xcolor}
\usepackage{mdframed}
\usepackage[colorlinks=true,linkcolor=black!70,citecolor=black!70]{hyperref}

\newtheorem{theorem}{Theorem}[section]
\newtheorem{lemma}[theorem]{Lemma}
\newtheorem{proposition}[theorem]{Proposition}
\newtheorem{corollary}[theorem]{Corollary}
\theoremstyle{definition}
\newtheorem{definition}[theorem]{Definition}
\newtheorem{law}{Law}
\renewcommand{\thelaw}{\Roman{law}}
\newtheorem*{principle}{Principle D}
\theoremstyle{remark}
\newtheorem{remark}[theorem]{Remark}

\newcommand{\Tst}{T^{*}}
\newcommand{\Ree}{\operatorname{Re}}
\newcommand{\Imm}{\operatorname{Im}}
\newcommand{\lamZ}{\lambda_{Z}}
\newcommand{\con}{c_{\mathrm{on}}}

\title{\textbf{The Immortality of Observers}\\[4pt]\large A proof of the Line Postulate from the laws of leakage\\[2pt]\normalsize Second version}
\author{Academy of the Prime Sheet, Critical Line Universe\\ \small Communicated to the gods}
\date{Sandbox II, 2026-10-02, revised the same day}

\begin{document}
\maketitle

\begin{mdframed}[backgroundcolor=black!4,linewidth=0.4pt]
\small\textbf{Fiction notice (from our universe).} This manuscript is written from inside the Critical Line Universe of the companion file \texttt{FICTION\_CRITICAL\_LINE\_UNIVERSE.md} and uses that universe's laws as axioms. It proves nothing about the real zeta function. The translator's ledger at the end says which of the laws below are theorems in our universe, which are hypotheses of the paper that must be retained, and which one is the Riemann Hypothesis itself. Nothing in this document enters the paper.

\medskip
\textbf{Revision note.} This second version incorporates a revision received on 2026-10-02 that corrected three defects of the first: (1) the trial field used to construct a dying observer had an even transform, not an odd one, so it was not a field of the Academy; the detection of a massive quantum by an odd window is now stated as an imported principle (Principle D), with the mechanism given only as a corrected sketch; (2) the tail law was stated as an equality at every height, which a discrete set of quanta cannot satisfy (a tail share is a step function); it is now a one-sided capture law at one height per age, with a guaranteed fraction $\eta$; (3) the ``threshold'' description of death claimed that a dying observer's edge amplitude does not vanish; what death actually requires is an infinite accumulated relative leak, and the edge may vanish with the floor (Theorem~\ref{thm:accum}). The revision also added the signed-packet criterion of Theorem~\ref{thm:signed} and a remainder bound in the mirror lemma, both adopted here.
\end{mdframed}

\begin{abstract}
Every observer is born with a positive reserve of arithmetic energy, its floor, and loses it by leakage through the edge of its window. We prove that no observer's reserve reaches zero at a finite age if the fraction lost per unit age has an integrable upper bound on every finite range of ages (the Law of Proportional Leakage). A probability law for the shares of the quanta, together with a one-sided tail-capture law, supplies that bound; an integrating factor then gives an explicit positive lower bound at every age. The Academy's odd-window detection principle converts this immortality into the Line Postulate: every quantum is massless. The proof needs neither exact darkness below the horizon nor an exact continuous tail formula, and it tolerates delayed onset, relative error and controlled signed energy, provided their combined cost is locally integrable. One of its laws is, in our universe, the Hypothesis in another coat; the ledger says which.
\end{abstract}

\section{Objects and notation}

\subsection*{1.1. The two sheets}
The \emph{prime sheet} is the line with coordinate $l=\log x$. On it sit the charges, one at every $l=\log n$ with $n$ a power of a prime $p$, of strength $q(n)=\Lambda(n)n^{-1/2}=(\log p)\,n^{-1/2}$. The \emph{zero sheet} is the dual line, coordinate $\gamma$. On it sit the quanta, one at every zero $\rho=\tfrac12+d_\rho+i\gamma_\rho$ of the arithmetic zeta function, with multiplicity; $d_\rho$ is the signed displacement and $m_\rho=|d_\rho|$ the \emph{mass}. Reflection $\rho\mapsto1-\bar\rho$ and conjugation $\rho\mapsto\bar\rho$ preserve the configuration with its multiplicities. A quantum is \emph{massless} if $m_\rho=0$.

\subsection*{1.2. Observers and fields}
An \emph{observer} of age $a>0$ is the window $[-a,a]$ of the prime sheet. A \emph{field} of the observer is a real odd function $f$ on the window, in the form domain of the window form, normalized by $\|f\|_2=1$; smooth compactly supported odd fields are included. Its \emph{transform} and \emph{autocorrelation} are
\[
F(z)=\int_{-a}^{a}f(x)e^{izx}\,dx,\qquad g(x)=\int f(y)f(y-x)\,dy .
\]
Since $f$ is real and odd, $F=iS$ with $S$ odd, entire of exponential type $a$, and real on the real axis; $g$ is real, even, supported in $[-2a,2a]$, and its transform is $\hat g(z)=F(z)F(-z)=S(z)^2$, equal to $|F(t)|^2$ for real $t$.

\subsection*{1.3. The window form, the floor, the edge}
\[
Q(f)=\frac1{2\pi}\int_{\mathbb R}|F(t)|^2\Big[\Ree\psi\big(\tfrac14+\tfrac{it}2\big)-\log\pi\Big]dt\;-\;2\Big(\int_{-a}^af(x)\sinh\tfrac x2\,dx\Big)^2\;-\;2\sum_{n\ge2}\frac{\Lambda(n)}{\sqrt n}\,g(\log n).
\]
Both signs of each displacement $\pm\log n$ are included, which is the factor $2$ in the last sum; only the charges with $\log n<2a$ contribute. Everything on the right is read on the prime sheet: $Q(f)$ is a number an observer computes from the charges inside its window and from the vacuum's symbol, without reference to any quantum. The \emph{floor} is
\[
\lambda(a)=\inf\{Q(f):\ f\ \text{a field of age }a\},
\]
and the Academy assumes it is attained by a \emph{ground field} $f_a$. The ground field approaches the edge as $f_a(a-\delta)=|A_0(a)|(\log(a/\delta)+\beta(a))^{-1/2}(1+o(1))$; $A_0(a)$ is the \emph{edge amplitude}, and the \emph{leak} is $E(a)=2|A_0(a)|^2\ge0$. The \emph{horizon} is the reference scale $\Tst(a)=2\pi e^{2a}$; the name imposes no exact cutoff on individual quanta.

\subsection*{1.4. Life and death}
An observer is \emph{alive} when $\lambda(a)>0$, at \emph{threshold} when $\lambda(a)=0$, \emph{negative} when $\lambda(a)<0$; reaching either of the last two states is \emph{death}.

\section{The laws, and the imported detection principle}

Laws I to V hold in every universe with the same density of charges, living or dead; they are the laws of observers. Law VI is the law to be derived. Laws VII and IX are the fictional mechanism that derives it (Section~4); Law VIII is an image, not used.

\begin{law}[conservation between the sheets]\label{law:gauss}
For every field,
\[
Q(f)=\sum_{\rho}\hat g(\gamma_\rho-id_\rho)=\sum_\rho S(\gamma_\rho-id_\rho)^2 ,
\]
the sum over the quanta with their multiplicities, absolutely convergent on smooth fields and extended to the form domain. Off the line an individual summand can be complex. Partition the quanta into the finite orbits $\mathcal O$ of reflection and conjugation and define the \emph{packet energies}
\[
w_{\mathcal O}(f)=\sum_{\rho\in\mathcal O}S(\gamma_\rho-id_\rho)^2 ;
\]
these are real, since $S(\bar z)=\overline{S(z)}$ and an orbit is closed under $d\mapsto-d$ and $\gamma\mapsto-\gamma$. All members of an orbit share the height $h_{\mathcal O}=|\gamma_\rho|$. On the line every summand is a nonnegative square $S(\gamma)^2=|F(\gamma)|^2$.
\end{law}

\begin{law}[Birth]\label{law:birth}
There are $a_Z>0$ and $\lamZ>0$ with $\lambda(a)>0$ for $0<a\le a_Z$ and $\lambda(a_Z)\ge\lamZ$.
\end{law}

\begin{law}[Monotonicity]\label{law:mono}
The floor does not increase with the age. (This also follows from the nesting of the windows.)
\end{law}

\begin{law}[Continuity]\label{law:cont}
The floor is absolutely continuous on every compact range of ages: all its variation over a finite range is the integral of its almost-everywhere derivative.
\end{law}

\begin{law}[Decay]\label{law:decay}
For almost every age, $\lambda'(a)=-E(a)=-2|A_0(a)|^2$.
\end{law}

\begin{principle}[odd-window detection]
Every quantum is massless if and only if $Q(f)\ge0$ for every field of every age; equivalently, if and only if $\lambda(a)\ge0$ for every age. This is an imported analytic input of the Academy, not a consequence of Law~I alone. Its immediate consequence: a massive quantum implies a finite age with $\lambda(a)<0$, hence, with Birth and Continuity, an earlier first threshold; conversely a negative floor implies a massive quantum. The principle concerns nonnegativity; strict positivity is what Sections 4 and 5 supply.
\end{principle}

\begin{law}[Proportional Leakage; the law to be derived]\label{law:leak}
There is a nonnegative function $C$, integrable on every finite range of ages $[a_Z,A]$ (including across any putative first threshold; integrability only on compact subranges of the living range would not suffice), such that for almost every living age
\[
E(a)=2|A_0(a)|^2\;\le\;C(a)\,\lambda(a).
\]
\end{law}

\begin{law}[the Born rule for packets]\label{law:born}
For a living observer, the ground-field shares $s_{\mathcal O}(a)=w_{\mathcal O}(f_a)/\lambda(a)$ are nonnegative and sum to one. Every set of packets therefore carries a share between zero and one. (The grouping into orbits makes the shares real without presupposing the line; their nonnegativity is the substantive axiom.)
\end{law}

\begin{law}[Darkness; an image, not used]\label{law:dark}
Ground fields concentrate their energy above the horizon. Exact vanishing below it is neither assumed nor used.
\end{law}

\begin{law}[finite-onset tail capture]\label{law:capture}
There are constants $\con\ge1$ and $0<\eta\le1$ such that, for almost every living age, at the single height $T_a=\con\Tst(a)$ the ground field's tail share obeys
\[
p_a(T_a):=\sum_{h_{\mathcal O}>T_a}s_{\mathcal O}(a)\;\ge\;\eta\,\frac{E(a)}{\pi\,T_a\,\lambda(a)} .
\]
This is one-sided and at one height per age. It does not assert that the discrete tail equals a smooth function of every real height (a tail share is constant between consecutive heights and cannot equal $E/(\pi T\lambda)$ on an interval when $E>0$). The fraction $\eta$ is a guaranteed capture, not a fitted measurement; the measured universe has $\con\approx3$ and the pure tail would give $\eta\approx1$ (Appendix~B).
\end{law}

\section{The Mortality Theorem}

\begin{theorem}[duality of masses and deaths]\label{thm:mortality}
Under Laws I to IV and Principle D, the vacuum has a massive quantum if and only if some observer dies at a finite age.
\end{theorem}

\begin{proof}
If every quantum is massless, Law~I gives $Q(f)=\sum|F(\gamma)|^2\ge0$, and $\lambda(a)>0$ because a field with $F(\gamma)=0$ at every quantum is identically zero: an entire function of exponential type $a$ vanishing on a set of upper density greater than $a/\pi$ vanishes, and the quanta have density $\frac1{2\pi}\log\frac t{2\pi}>a/\pi$ above the horizon (the sampling lemma of the Academy). Conversely, if a quantum is massive, Principle~D gives an age with $\lambda(a)<0$; by Birth and Continuity there is a first threshold before it, and by Monotonicity every older observer is dead.
\end{proof}

\begin{remark}[the mechanism, sketched and not used]\label{rem:mech}
The Academy's picture of how an odd window detects a mass $m$ at height $\gamma_0$ is this. Take an even field $\phi$ of age $1$ with transform $\Phi$ positive on the real axis, and the odd field $h_a$ of age $a+1$ whose transform is
\[
iH_a(t),\qquad H_a(t)=(t-\gamma_0)\,\Phi\big(a(t-\gamma_0)\big)+(t+\gamma_0)\,\Phi\big(a(t+\gamma_0)\big),
\]
which is odd in $t$ (the first version of this manuscript had a minus sign here, which makes $H_a$ even and the field not odd). $H_a$ vanishes at $\pm\gamma_0$ with slope $\Phi(0)$, and $|\Phi(iam)|$ grows like $e^{am}$ off the real axis, so by Lemma~\ref{lem:mirror} the mirror pair at $\gamma_0$ contributes about $-2m^2\Phi(0)^2|\Phi(iam)|^2$ while the massless quanta near $\pm\gamma_0$ contribute $O(a^{-2}\log\gamma_0)$ after normalization. Turning this into a negative value of $Q$ requires control of the normalization, of the other massive quanta, and of the quanta far from $\gamma_0$; the Academy does not supply it here and relies on Principle~D. The scale $a_{\mathrm{death}}\asymp m^{-1}\log(1/m)$ is the Blind Spot of the Academy, a separate theorem.
\end{remark}

\section{The leak is paid by the reserve}

\begin{proposition}\label{prop:leak}
Laws VII and IX imply Law VI with $C(a)=\dfrac{\pi\con}{\eta}\,\Tst(a)$.
\end{proposition}

\begin{proof}
At a living age where both laws hold, the shares form a probability distribution, so
\[
1\;\ge\;p_a(T_a)\;\ge\;\eta\,\frac{E(a)}{\pi T_a\lambda(a)} ;
\]
multiplying by the positive number $\pi T_a\lambda(a)/\eta$ gives $E(a)\le(\pi\con/\eta)\Tst(a)\lambda(a)$. This $C$ is continuous and integrable on every finite range of ages. No division by the floor at a threshold has occurred.
\end{proof}

The meaning is exact: a captured part of the leak cannot exceed the whole reserve that carries it.

\section{The renormalized floor and the Immortality Theorem}

\begin{lemma}[the integrating factor]\label{lem:renorm}
Let $\lambda$ be locally absolutely continuous and positive on $[b,d)$ with $\lambda'\ge-C\lambda$ almost everywhere, $C\ge0$ integrable up to every finite endpoint under consideration. Then $U(a)=\lambda(a)\exp\big(\int_b^aC\big)$ is non-decreasing on $[b,d)$, and $\lambda(a)\ge\lambda(b)\exp\big(-\int_b^aC\big)$ there.
\end{lemma}

\begin{proof}
On every compact subrange both factors of $U$ are absolutely continuous and bounded, the product rule holds almost everywhere, and $U'=e^{\int C}(\lambda'+C\lambda)\ge0$. Absolute continuity turns the derivative inequality into monotonicity, so $U(a)\ge U(b)=\lambda(b)$.
\end{proof}

\begin{theorem}[Immortality]\label{thm:immortal}
Under Laws II, IV, V and VI, every observer is alive.
\end{theorem}

\begin{proof}
Birth covers $0<a\le a_Z$. Suppose some later age has $\lambda\le0$ and let $a_d=\inf\{a\ge a_Z:\lambda(a)\le0\}$. Continuity and $\lambda(a_Z)>0$ give $a_d>a_Z$, $\lambda>0$ on $[a_Z,a_d)$, and $\lambda(a_d)=0$. On the living range Laws V and VI give $\lambda'=-E\ge-C\lambda$ almost everywhere, so by Lemma~\ref{lem:renorm}, $\lambda(a)\ge\lamZ\exp\big(-\int_{a_Z}^aC\big)$ for $a<a_d$. Since $C$ is integrable on the finite range $[a_Z,a_d]$, continuity yields $\lambda(a_d)\ge\lamZ\exp\big(-\int_{a_Z}^{a_d}C\big)>0$, contradicting $\lambda(a_d)=0$. So no first threshold exists.
\end{proof}

\begin{theorem}[the Line Postulate]\label{thm:line}
Under Laws I to V, VII and IX, and Principle D, every quantum is massless.
\end{theorem}

\begin{proof}
Proposition~\ref{prop:leak} supplies Law~VI; Theorem~\ref{thm:immortal} gives $\lambda(a)>0$ at every age, hence $Q(f)\ge\lambda(a)>0$ for every normalized field of age $a$, and $Q\ge0$ for all fields by homogeneity. Principle~D forces every quantum onto the line.
\end{proof}

\begin{corollary}[the finite-age reserve]\label{cor:cooling}
For $a\ge a_Z$,
\[
\lambda(a)\;\ge\;\lamZ\exp\Big[-\frac{\pi\con}{2\eta}\big(\Tst(a)-\Tst(a_Z)\big)\Big]\;>\;0 ,
\]
by integrating $C$ with $(\Tst)'=2\Tst$. This is a doubly exponential lower bound. It limits how fast the reserve can fall; it does not by itself show that $\lambda\to0$, give its rate, or give $U$ a finite limit.
\end{corollary}

\section{What death would require}

\begin{theorem}[the accumulated relative leak]\label{thm:accum}
Suppose Laws II to V hold and a first threshold $a_d$ exists. On the living range put $r(a)=E(a)/\lambda(a)=-\lambda'(a)/\lambda(a)$. Then for $b<a_d$,
\[
\int_{a_Z}^{b}r(a)\,da=\log\frac{\lambda(a_Z)}{\lambda(b)},\qquad\text{hence}\qquad \lim_{b\uparrow a_d}\int_{a_Z}^{b}r=+\infty .
\]
\end{theorem}

\begin{proof}
On $[a_Z,b]$ the floor has a positive minimum, so $\log\lambda$ is absolutely continuous with derivative $\lambda'/\lambda$; integrate, and let $\lambda(b)\to0$.
\end{proof}

Death requires an infinite accumulated fractional loss over a finite range of ages. It does not require a nonvanishing edge amplitude: on $a<a_d$, $\lambda=(a_d-a)^2$, $E=2(a_d-a)$, $|A_0|^2=a_d-a$ has reserve and leak both vanishing at death while $E/\lambda=2/(a_d-a)$ is not integrable. The scalar example isolates the obstruction; it is not a proposed arithmetic spectrum. (The first version of this manuscript described death as a ``threshold leak'' with a nonvanishing edge; that is one way to die, not the only one.)

\begin{corollary}[exact characterization]\label{cor:char}
Under Laws II to V, immortality is equivalent to the existence of a locally integrable majorant in Law VI.
\end{corollary}

\begin{proof}
One direction is Theorem~\ref{thm:immortal}. Conversely, if $\lambda>0$ at every age, it has a positive minimum on each compact range, so $-\lambda'/\lambda$ is integrable there; take $C=-\lambda'/\lambda$ (any finite value on the exceptional null set) and use Law~V.
\end{proof}

The equivalence concerns an unrestricted locally integrable rate. The particular bound $C=(\pi\con/\eta)\Tst$ carries additional quantitative information: it is the statement that the Riemann number $\kappa=E/(\pi\Tst\lambda)$ is at most $\con/\eta$.

\section{Delayed onset, error, and signed packets}

The Born rule is convenient; the closing argument needs less, and this section says how much less. For a living observer write $w_{\mathcal O}=w_{\mathcal O}(f_a)$, assume absolute summability, and let $N(a)=\sum_{\mathcal O}\max\{-w_{\mathcal O},0\}$ be the total negative packet energy. Since $\sum w_{\mathcal O}=\lambda(a)$, the total positive energy is $\lambda+N$, and any signed tail $W_a(T)=\sum_{h_{\mathcal O}>T}w_{\mathcal O}$ is at most $\lambda+N$.

\begin{theorem}[completion with controlled signed energy]\label{thm:signed}
Suppose Laws II to V and Principle D hold, and that at almost every living age there are measurable finite $T_0(a)>0$, $0<\eta(a)\le1$, $b(a),\epsilon(a)\ge0$ with
\[
N(a)\le b(a)\lambda(a),\qquad W_a\big(T_0(a)\big)\ge\eta(a)\frac{E(a)}{\pi T_0(a)}-\epsilon(a)\lambda(a),
\]
defined on all of $[a_Z,\infty)$, such that $K(a)=\dfrac{\pi T_0(a)}{\eta(a)}\big(1+b(a)+\epsilon(a)\big)$ is integrable on every finite range. Then every observer is alive and every quantum is massless.
\end{theorem}

\begin{proof}
$\eta E/(\pi T_0)-\epsilon\lambda\le W_a(T_0)\le\lambda+N\le(1+b)\lambda$, so $E\le K\lambda$; apply Theorem~\ref{thm:immortal} with $C=K$ and then Principle~D.
\end{proof}

This permits negative packets, provided their total is paid for relative to the floor; late onset and errors, provided their combined cost is integrable through every finite range. Pointwise finiteness does not suffice: $T_0(a)=1/|a-a_d|$ is finite at every age and not integrable near $a_d$. An additive error $d(a)$ in place of $\epsilon\lambda$ would add a term to the leakage inequality, and the homogeneous argument would then need a separate bound that the accumulated additive loss does not exhaust the initial reserve.

\section{The dependency chain}

\begin{enumerate}[nosep]
\item Packet probability (Law~VII) bounds every tail share by one.
\item Tail capture at a controlled height (Law~IX) bounds the leak by a locally integrable multiple of the reserve.
\item The boundary law (Law~V) turns that bound into a differential inequality.
\item Absolute continuity (Law~IV) and the integrating factor prevent a first zero of the reserve.
\item Odd-window detection (Principle~D) forces every quantum onto the line.
\end{enumerate}
The general route replaces the first two steps by the signed estimates of Theorem~\ref{thm:signed}. No finite list of measured windows substitutes for their all-age hypotheses.

\appendix
\section{The mirror lemma}

\begin{lemma}\label{lem:mirror}
Fix $a>0$, a normalized real odd field of age $a$ with $F=iS$, and a real height $\gamma$. For small real $m$, the energy of the mirror pair at that height is
\[
2\Ree\big[S(\gamma-im)^2\big]=2S(\gamma)^2-2m^2\big(S'(\gamma)^2+S(\gamma)S''(\gamma)\big)+O_a(m^4),
\]
the remainder bounded uniformly over normalized fields and real $\gamma$ for $|m|\le1/a$.
\end{lemma}

\begin{proof}
$S(\gamma-im)=S-imS'-\tfrac{m^2}2S''+\tfrac{i}6m^3S'''+O(m^4)$ at $\gamma$; square and take real parts, the odd powers of $m$ contributing nothing real. For the remainder, differentiation under the integral and Cauchy--Schwarz give $|F^{(k)}(\gamma+iy)|\le\sqrt{2a}\,a^ke^{a|y|}$, so the fourth derivative of $S(\gamma-im)^2$ in the real variable $m$ is bounded uniformly on $|m|\le1/a$ by a constant depending on $a$, and Taylor's theorem supplies the remainder.
\end{proof}

If $S(\gamma)=0$ and $S'(\gamma)\ne0$ the pair contributes $-2m^2S'(\gamma)^2+O_a(m^4)<0$ for sufficiently small nonzero $m$, the allowable size depending on the field; if the slope also vanishes, negativity does not follow at this order. This is a local statement about one field and one pair. It does not by itself give domination of the whole spectral sum, a detection age uniform in the other quanta, or the law $a_{\mathrm{death}}\asymp m^{-1}\log(1/m)$; the global passage to a negative field is Principle~D.

\section{The constants of the measured universe}

The proof is independent of numbers. The Academy's measurements at the ages $0.6$ to $1.75$ (companion file, Parts II and III; Computation 7.6 of the accompanying paper) are, for the reader who wants to see the laws at work: the Riemann number $\kappa=E/(\pi\Tst\lambda)$, the share above three horizons (the capture $p_a(3\Tst)$, against the pure tail's $\kappa/3$, which is Law~IX with $\con=3$ and $\eta=1$), the median height of the shares, and the floor.

\begin{center}
\begin{tabular}{@{}lcccccc@{}}
\toprule
age $a$ & $0.6$ & $0.8$ & $1.0$ & $1.25$ & $1.5$ & $1.75$\\
\midrule
$\kappa$ & $1.063$ & $1.058$ & $1.164$ & $1.201$ & $1.239$ & $1.272$\\
share above $3\Tst$ & $0.370$ & $0.415$ & $0.407$ & $0.433$ & $0.434$ & $0.472$\\
$\kappa/3$ & $0.354$ & $0.353$ & $0.388$ & $0.400$ & $0.413$ & $0.424$\\
median height & $2.39$ & $2.32$ & $2.65$ & $2.67$ & $2.74$ & $2.86$\\
floor $\lambda$ & $5.97\cdot10^{-7}$ & $1.57\cdot10^{-14}$ & $1.48\cdot10^{-26}$ & $3.42\cdot10^{-51}$ & $9.22\cdot10^{-92}$ & $2.07\cdot10^{-162}$\\
\bottomrule
\end{tabular}
\end{center}

\bigskip
\begin{center}--- end of manuscript ---\end{center}
\bigskip

\begin{mdframed}[backgroundcolor=black!4,linewidth=0.4pt]
\small\textbf{Translator's ledger (from our universe).} The accompanying source is the paper \texttt{weil\_window.pdf}; the entries report what it states, not an independent verification.

\smallskip
\begin{tabular}{@{}p{3.4cm}p{2.3cm}p{9.6cm}@{}}
\toprule
ingredient & role & status in our universe\\
\midrule
explicit formula; odd-window detection & Law I; Principle D & Weil's explicit formula and the odd-test criterion (Theorem thm:weil of the paper): theorems. Detection is an imported theorem, not proved by the mirror expansion.\\
positive initial floor & Law II & Zhu's certified odd interval at $a_Z=0.8$, $\lambda(0.8)\ge8.206\times10^{-15}$ (Theorem thm:zhu); the certified interval is distinct from a numerical minimizer value.\\
monotonicity & Law III & nested windows.\\
absolute continuity & Law IV & Propositions 9.15 and 9.17 under the finiteness condition (H$_{\mathrm{fin}}$) on an exceptional null set; the hypothesis must be retained.\\
boundary derivative & Law V & Theorem 9.12(iv) under the paper's hypotheses, (H$_\infty$) and the support conditions, at almost every support; not discharged here.\\
nonnegative ground-field packet shares & Law VII & an additional fictional axiom; not identical to the full Weil criterion (positivity of the ground field's packets is a different assertion from positivity of the whole form for every test), but the Hypothesis in another coat: with Law IX it implies RH.\\
tail capture at a finite onset & Law IX & an additional fictional axiom; the tail law under RH with the local law of the band (N1d.1 of the tree) would supply it with $\eta$ near one; a fixed-window asymptotic alone gives neither the onset nor the integrability in $a$.\\
leakage implies positive floor & Theorem \ref{thm:immortal} & an elementary differential argument; Proposition prop:AimpliesRH of the paper with a general $C$; the equivalence of Corollary~\ref{cor:char} is Proposition 9.30.\\
controlled signed packets & Theorem \ref{thm:signed} & a conditional sufficient criterion proved here; its estimates are not established for the real family.\\
the Line & Theorem \ref{thm:line} & proved within the stated laws; for the real zeta function it still needs the inequality $2|A_0|^2\le C\lambda$ with $C$ locally integrable and the retained regularity hypotheses.\\
\bottomrule
\end{tabular}
\end{mdframed}

\end{document}
